\subsection{Installation Filius}
\label{subsec:filius-installation}
\subsubsection{Windows}
Laden Sie folgende Datei herunter und installieren Sie das Programm:  \href{https://informatik.mygymer.ch/base/?page=network/2-filius/windows/}{Link}.

Falls Sie dazu aufgefordert werden, klicken Sie immer auf ``erlauben'', bzw. ``weiter''. Nachdem Sie das Programm installiert haben, können Sie folgende Schritte ausführen:
\begin{enumerate}
	\item Das heruntergeladene Installationsprogramm (mit der Dateiendung \texttt{.exe}) löschen
	\item Auf der Tastatur tippen: \faWindows+\texttt{Filius}
	\item Filius öffnen
\end{enumerate}
\subsubsection{MacOS}
\begin{mycaution}
	Falls Probleme beim Ausführen der unten stehenden Schritte auftreten, schauen Sie sich die genauen Anleitungen mit Screenshots auf \href{https://oinf.ch/wp-content/uploads/Anleitung_Filius_Mac-OS.pdf}{dieser Webseite} an. Sie können natürlich auch mich um Hilfe fragen \emoji{winking-face}.
\end{mycaution}
\paragraph*{Installation Java SDK (Voraussetzung für Filius)}
Damit Filius ausgeführt werden kann, muss zuerst Java SDK in der neusten Version installiert werden.
\begin{enumerate}
	\item Laden Sie Java als DMG-Installer unter folgendem Link herunter: \url{https://www.oracle.com/java/technologies/downloads/}
	      \begin{itemize}
		      \item Überprüfen Sie, ob Sie über einen Mac mit Apple- oder Intel-Chip verfügen, indem Sie ganz oben links auf Ihrem Bildschirm auf das \faApple{}-Zeichen $\rightarrow$ ``Über diesen Mac'' klicken und die Beschreibung Ihres Chips (erste Zeile) lesen.
		      \item Falls Sie einen neueren Mac mit Apple-Silicon-Chip (z.B. M1, M2, M3...) haben (Modelle nach 2022 gekauft), laden Sie die ``ARM''-Version herunter (als \texttt{.dmg}-Datei)
		      \item Falls Sie einen Mac mit Intel-Chip haben, laden Sie die ``x64''-Version herunter
	      \end{itemize}
\end{enumerate}

Installieren Sie Java SDK, indem Sie auf die heruntergeladene Datei klicken.

\paragraph*{Installation Filius}

\begin{itemize}
	\item Laden Sie Filius in der neusten Version \textbf{als \texttt{.zip}-Datei} von folgendem Link herunter: \url{https://www.lernsoftware-filius.de/Herunterladen}.
	\item Verschieben (ziehen) Sie den Ordner in Ihren Anwendungs-Ordner.
	\item Öffnen Sie den Ordner und starten Sie die Datei \texttt{filius.jar}, indem Sie rechts darauf klicken (\keys{\Alt} + Klick) und ``Öffnen'' wählen. Dies müssen Sie nur beim ersten Mal machen, danach können Sie ``normal'' auf das Programm klicken (ohne \keys{\Alt}), um es zu öffnen.
\end{itemize}